\documentclass[../main.tex]{subfiles}

\begin{document}

\ifSubfilesClassLoaded{

    %\tableofcontents
    
    \setcounter{chapter}{3}
}{}

\chapter{ Week 4 }
代靖涵 25120222201319
\begin{exercise}{}{}
    某班级学生的考试成绩数学不及格的占 8\%, 语文不及格的占 5\%, 这两门都不及格的占 2\%.
\begin{enumerate}
    \item 已知一学生数学不及格, 他语文也不及格的概率是多少?
    \item 已知一学生语文不及格, 他数学也不及格的概率是多少?
\end{enumerate}
\end{exercise}

\begin{proof}
    设数学, 语文, 不及格的事件分别为\(  A  \), \(  B  \). 则 \[
    P\left( A \right)= 0.08,\quad P\left( B \right)= 0 . 0 5  ,\quad P\left( AB \right)= 0.02 
    \]
    \begin{enumerate}
        \item 代表的事件为 \( B|A    \), 概率为 \[
        \begin{aligned}
        P\left( B | A \right)= \frac{P\left( AB \right)  }{P\left( A \right)  }= 0.25   
        \end{aligned}
        \]
        \item 代表的事件为 \(  A|B  \), 概率为 \[
        P\left( A|B \right)= \frac{P\left( AB \right)  }{P\left( B \right)  }= 0.4  
        \]  
    \end{enumerate}
      
\end{proof}
\begin{exercise}{}{}
掷两颗骰子, 以 $A$ 记事件“两颗点数之和为 10”, 以 $B$ 记事件“第一颗点数小于第二颗点数”, 试求条件概率 $P(A|B)$ 和 $P(B|A)$.
\end{exercise}

\begin{proof}
    考虑试验的样本空间为 \(  \Omega =  \left\{ 1,2,3,4,5,6 \right\}^{2}  \), 则 \(  \left|  \Omega  \right|= 36   \). \(  \left| A \right|= 3  \),  \[
    P\left( A \right)= \frac{\left| A \right|  }{\left|  \Omega  \right|  }= \frac{3 }{36 }   = \frac{1 }{12 } 
    \]  \(  \left| B \right|= 5+ 4+ 3+ 2+ 1= 15   \), \[
    P\left( B \right)= \frac{\left| B \right|  }{\left|  \Omega  \right|  }= \frac{15 }{36 }= \frac{5 }{12 }    
    \]  \(  \left| AB \right|=  1  \), \[
    P\left( AB \right)= \frac{\left| AB \right|  }{\left|  \Omega  \right|  }= \frac{1 }{36 }   
    \]因此 \[
    P\left( A|B \right)= \frac{P\left( AB \right)  }{P\left( B \right)  }= \frac{1 }{15 }   
    \] \[
    P\left( B|A \right)= \frac{P\left( AB \right)  }{P\left( A \right)  }= \frac{1 }{3 }   
    \] 
\end{proof}
\begin{exercise}{}{}
设 10 件产品中有 3 件不合格品,从中任取两件,已知其中一件是不合格品,求另一件也是不合格品的概率.
\end{exercise}

\begin{proof}
设至少有一件不合格的事件为 \(  A  \), 则 \[
P\left( A \right)= 1-P\left( \bar{A} \right)= 1-\frac{C_{7}^{2} }{C_{10}^{2} }=\frac{8 }{15 } 
\] 设两件均不合格的事件为 \(  B  \), 则 \[
P\left( B \right)= \frac{C_{3}^{2} }{C_{10}^{2} }= \frac{1 }{15 }   
\]则已知其中一件不合格, 另一件也不合格的事件为 \(  B|A  \). 此外注意到 \(  B\subseteq A  \), 因此 \(  B|A  \)的  概率为 \[
P\left( B|A \right)= \frac{P\left( A\cap B \right)  }{P\left( A \right)  }= \frac{P\left( B \right)  }{ P\left( A \right) }= \frac{1 }{8 }    
\]
\end{proof}
\begin{exercise}{}{}
已知 $P(A) = 1/3, P(B|A) = 1/4, P(A|B) = 1/6$, 求 $P(A \cup B)$.
\end{exercise}

\begin{proof}
     \[
     P\left( A\cup B \right)= P\left( A \right)+ P\left( B \right)-P\left( A\cap B \right)    
     \]其中 \[
     P\left( A\cap B \right)= P\left( A \right)P\left( B|A \right)= \frac{1 }{12 }    
     \] \[
     P\left( B \right)= \frac{P\left( A\cap B \right)  }{P\left( A|B \right)  }  = \frac{1 }{2 } 
     \]因此 \[
     P\left( A\cup B \right)= \frac{1 }{3 }+ \frac{1 }{2 }-\frac{1 }{12 }= \frac{3 }{4 }     
     \]
\end{proof}
\begin{exercise}{}{}
已知 $P(\bar{A}) = 0.3, P(B) = 0.4, P(A\bar{B}) = 0.5$, 求 $P(B|A \cup \bar{B})$.
\end{exercise}

\begin{proof}
    \[
    P\left( B|A\cup \bar{B} \right)= \frac{P\left( B\cap \left( A\cup \bar{B} \right)  \right)  }{P\left( A\cup \bar{B} \right)  }  
    \]其中 \[
    \begin{aligned}
    P\left( B\cap \left( A\cup \bar{B} \right)  \right)&= P\left( AB\cup B \bar{B} \right)= P\left( AB \right)    
    \end{aligned}
    \]又注意到 \[
    P\left( AB \right)= P\left( A \right)-P\left( A \bar{B} \right)   = 1-P\left( \bar{A} \right)-P\left( A \bar{B} \right)= 0.2  
    \]因此 \[
    P\left( B\cap \left( A\cup \bar{B} \right)  \right)= 0.2 
    \]此外 \[
   \begin{aligned}
    P\left( A\cup \bar{B} \right)&= P\left( A \right)+ P\left( \bar{B} \right)-P\left(  A \bar{B} \right)    \\ 
     &= 1-P\left( \bar{A} \right)+ 1-P\left( B \right)-P\left( A \bar{B} \right)\\ 
      &= 0.8   
   \end{aligned} 
    \]因此 \[
    P\left( B|A\cup \bar{B} \right)= \frac{1 }{4 }  
    \]
\end{proof}



\begin{exercise}{}{}
设 $A, B$ 为两事件, $P(A) = P(B) = 1/3, P(A|B) = 1/6$, 求 $P(A|\bar{B})$.
\end{exercise}

\begin{proof}
    \[
    P\left( A|\bar{B} \right)= \frac{P\left( A \cap \bar{B} \right)  }{P\left( \bar{B} \right)  }  
    \]其中 \[
   \begin{aligned}
    P\left( A\cap \bar{B} \right)&= P\left( A \right) -P\left( A\cap B \right)\\ 
     &= P\left( A \right) - P\left( A|B \right)P\left( B \right)\\ 
      &= \frac{5 }{18 }     
   \end{aligned}
    \] \[
    P\left( \bar{B} \right)= 1-P\left( B \right)= \frac{2 }{3 }   
    \]因此 \[
    P\left( A|\bar{B} \right)= \frac{\frac{5 }{18 }  }{\frac{2 }{3 }  }  = \frac{5 }{12 } 
    \]
\end{proof}
\begin{exercise}{}{}
口袋中有 1 个白球, 1 个黑球. 从中任取 1 个, 若取出白球, 则试验停止; 若取出黑球, 则把取出的黑球放回的同时, 再加入 1 个黑球, 如此下去, 直到取出的是白球为止, 试求下列事件的概率:
\begin{enumerate}
    \item 取到第 $n$ 次, 试验没有结束;
    \item 取到第 $n$ 次, 试验恰好结束.
\end{enumerate}
\end{exercise}

\begin{proof}
     \begin{enumerate}
        \item  令 $B_k$ 为事件 "第 $k$ 次取出的是黑球".   则 \[
   \begin{aligned}
        P(A)&= P\left( B_1\cap B_2\cap \cdots \cap B_{n} \right) \\ 
         & = P(B_1) P(B_2|B_1) P(B_3|B_1 \cap B_2) \cdots P(B_n|B_1 \cap B_2 \cap \cdots \cap B_{n-1})  \\ 
          &= \frac{1 }{2 }\cdot \frac{2 }{3 }\cdot \cdots \cdot \frac{n }{n+ 1 }\\ 
           &= \frac{1 }{n+ 1 }    
   \end{aligned}
        \] 
        \item 记取到第 \(  n-1  \)次, 试验没有结束的事件为 \(  B  \), 则 \[
        P\left( B \right)= \frac{1 }{n }  
        \]  则取到第 \(  n  \)次, 试验恰好结束, 可以表述为``事件在取到第 \(  n  \)次时已经结束,并且 在取到第 \(  n-1  \)次时没有结束  '', 即事件 \(  \bar{A}B \), 又注意到 \( A\subseteq B\), 因此 概率\[
        P\left( \bar{A}B \right)= P\left( B \right)-P\left( AB \right)= P\left( B \right)-P\left( A \right)= \frac{1 }{n\left( n+ 1 \right)  }      
        \]
     \end{enumerate}
     
\end{proof}

\begin{exercise}{}{}
将 $n$ 根绳子的 $2n$ 个头任意两两相接, 求恰好结成 $n$ 个圈的概率.
\end{exercise}

\begin{proof}
  令 \(  A_{i}  \)表示第 \(  i  \)根绳子自成一圈的概率, 则恰好结成 \(  n  \)个圈的概率为 \[
 \begin{aligned}
  P\left( A_1A_2\cdots A_{n} \right)&= P\left( A_1 \right)P\left( A_2|A_1 \right)\cdots P\left( A_{n} |A_1A_2\cdots A_{n-1}\right)    \\ 
   &= \frac{1 }{2n-1 }\cdot \frac{1 }{2\left( n-1 \right)-1  }\cdots \frac{1 }{1 }\\ 
    &= \frac{1 }{\left( 2n-1 \right)!!  }    
 \end{aligned}
  \]   
\end{proof}

\begin{exercise}{}{}
    设 $P(A)>0$, 证明:
$$P(B|A) \ge 1 - \frac{P(\bar{B})}{P(A)}.$$
\end{exercise}

\begin{proof}
     \[
     P\left( A \right)P\left( B|A \right)+ P\left( \bar{B} \right)= P\left( AB \right)+ P\left( \bar{B} \right)\ge P\left( AB \right)+ P\left( A \bar{B} \right)= P\left( A \right)        
     \]因此 \[
     P\left( A \right)P\left( B|A \right)+ P\left( \bar{B} \right)\ge P\left( A \right)    
     \]两边除以 \(  P\left( A \right)   \)并移项, 即得原不等式成立. 
\end{proof}

\begin{exercise}{}{}
    若事件 $A$ 与 $B$ 互不相容, 且 $P(\bar{B}) \ne 0$, 证明:
$$P(A|\bar{B}) = \frac{P(A)}{1-P(B)}.$$
\end{exercise}

\begin{proof}
    若事件 \(  A  \)与 \(  B  \)互不相容, 则 \(  P\left( AB \right)= 0   \), 我们有 \[
    P\left( A|\bar{B} \right)= \frac{P\left( A \bar{B} \right)  }{P\left( \bar{B} \right)  }= \frac{P\left( A \right)-P\left( AB \right)   }{1-P\left( B \right)  }= \frac{P\left( A \right)  }{1-P\left( B \right)  }    
    \]   
\end{proof}

\begin{exercise}{}{}
    设 $A, B$ 为任意两个事件, 且 $A \subset B, P(B)>0$, 证明 $P(A) \le P(A|B)$.
\end{exercise}

\begin{proof}
    若 \(  A\subseteq B  \), 则 \(  P\left( AB \right)= P\left( A \right)    \), 且利用 \(  P\left( B \right)\le 1   \),  我们得到 \[
    P\left( A|B \right)= \frac{P\left( AB \right)  }{P\left( B \right)  }= \frac{P\left( A \right)  }{P\left( B \right)  }  \ge P\left( A \right)  
    \]  
\end{proof}
\begin{exercise}{}{}
若 $P(A|B)>P(A|\bar{B})$,证明 $P(B|A)>P(B|\bar{A})$.
\end{exercise}

\begin{proof}
    注意到 \(  P\left( A|B \right)> P\left( A| \bar{B} \right)    \)当且仅当 \[
    \frac{P\left(  AB \right)  }{P\left( B \right)  }> \frac{P\left( A \bar{B} \right)  }{P\left( \bar{B} \right)  } \iff  \frac{P\left( AB \right)  }{P\left( B \right)  }> \frac{P\left( A \right)-P\left( AB \right)   }{1-P\left( B \right)  }\iff P\left( AB \right)> P\left( A \right)P\left( B \right)      
    \]通过调换 \(  A,B  \)的位置, 可知 \(  P\left( B|A \right)> P\left( B|\bar{A} \right)    \)也当且仅当 \[
    P\left( AB \right)> P\left( A \right)P\left( B \right)   
    \]  因此原不等式成立.
\end{proof}

\begin{exercise}{}{}
设 $P(A)=p, P(B)=1-\varepsilon$,证明:
\[ \frac{p-\varepsilon}{1-\varepsilon} \le P(A|B) \le \frac{p}{1-\varepsilon}. \]
\end{exercise}

\begin{proof}
    \[
   \begin{aligned}
    P\left( A|B \right)= \frac{P\left( AB \right)  }{P\left( B \right)  }&= \frac{P\left( A \right)+ P\left( B \right)-P\left( A\cup B \right)      }{ P\left( B \right) }\\ 
     &\ge \frac{P\left( A \right)+ P\left( B \right)-1  }{P\left( B \right)  }\\ 
      &= \frac{p+ 1- \varepsilon -1  }{1- \varepsilon  }\\ 
       &= \frac{p- \varepsilon  }{1- \varepsilon  }  
   \end{aligned}    
    \]  \[
    P\left( A|B \right)= \frac{P\left( AB \right)  }{P\left( B \right)  }\le \frac{P\left( A \right)  }{P\left( B \right)  }= \frac{p }{1- \varepsilon  }    
    \]
\end{proof}

\begin{exercise}{}{}
若 $P(A|B)=1$,证明 $P(\bar{B}|\bar{A})=1$.
\end{exercise}
\begin{proof}
   注意到  \[
     1= P\left( A|B \right)= \frac{P\left( AB \right)  }{P\left( B \right)  }\implies P\left( B \right)= P\left( AB \right)     
     \]于是 \[
     P\left( B|\bar{A} \right)= \frac{P\left( \bar{A}B \right)  }{P\left( \bar{A} \right)  }= \frac{P\left( B \right)-P\left( AB \right)   }{P\left( \bar{A} \right)  }= 0   
     \]因此 \[
     P\left( \bar{B}|\bar{A} \right) =1- P\left( B|\bar{A} \right) = 1
     \]
\end{proof}

\begin{exercise}{}{}
一盒晶体管中有 8 只合格品、2 只不合格品.从中不返回地一只一只取出,试求第二次取出的是合格品的概率.
\end{exercise}

\begin{proof}
    记第 \(  i  \)次取出合格品的事件为 \(  A_{i}  \), \(  i= 1,2  \). 则由全概率公式 \[
    \begin{aligned}
    P\left( A_{2} \right)&=  P\left( A_{2}|A_{1} \right)P\left( A_{1} \right)+ P\left( A_{2}|A_{1}^{c} \right)P\left( A_{1}^{c} \right)\\ 
     &= \frac{7 }{9 }\cdot \frac{4 }{5 }+ \frac{8 }{9 }\cdot \frac{1 }{5 }\\ 
      &= \frac{4 }{5 }           
    \end{aligned}
    \]   
\end{proof}
\begin{exercise}{}{}
甲口袋有 $a$ 个白球、$b$ 个黑球,乙口袋有 $n$ 个白球、$m$ 个黑球.
\begin{enumerate}
    \item 从甲口袋任取 1 个球放入乙口袋,再从乙口袋任取 1 个球.试求最后从乙口袋取出的是白球的概率;
    \item 从甲口袋任取 2 个球放入乙口袋,再从乙口袋任取 1 个球.试求最后从乙口袋取出的是白球的概率.
\end{enumerate}
\end{exercise}

\begin{proof}
    \begin{enumerate}
        \item 设第一次从甲口袋中取出白球的事件为 \(  B_{W}  \), 取出黑球的事件为 \(  B_{B}  \)  , 则 \(  B_{W}\cap B_{B}= \varnothing, B_{W}\cup B_{B}=  \Omega   \). 设最后从乙口袋取出白球的事件为 \(  A_{W}  \), 由全概率公式, 我们有 \[
        \begin{aligned}
        P\left( A_{W} \right)&= P\left( A_{W}|B_{B} \right)P\left( B_{B} \right)+ P\left( A_{W}|B_{W} \right)P\left( B_{W} \right)\\ 
         &= \frac{n }{m+ n+ 1 }\cdot \frac{b }{a+ b }        + \frac{ n+ 1}{m+ n+ 1 } \cdot \frac{a }{a+ b } \\ 
          &= \frac{nb+ \left( n+ 1 \right)a  }{ \left( m+ n+ 1 \right)\left( a+ b \right)  } 
        \end{aligned}
        \]  
        \item 设从甲口袋中取出两个白球的事件为 \(  B_{WW}  \), 取出两个黑球的事件为 \(  B_{BB}  \), 取出一黑一白的事件为 \(  B_{BW}  \). 则 \[
        P\left( B_{WW} \right)= \frac{C_{a}^{2} }{C_{a+ b}^{2} },\quad P\left( B_{WB} \right)= \frac{C_{a}^{1}C_{b}^{1} }{C_{a+ b}^{2} }    ,\quad P\left( B_{BB} \right)= \frac{C_{b}^{2} }{C_{a+ b}^{2} }  
        \]   设最后从乙口袋取出白球的事件为 \(  A_{W}  \),  由全概率公式, 我们有 \[
        \begin{aligned}
        P\left( A_{W} \right)&= P\left( A_{W}|B_{WW} \right)P\left( B_{WW} \right)+ P\left( A_{W}|B_{WB} \right)P\left( B_{WB} \right)+ P\left( A_{W}|B_{BB} \right)P\left( B_{BB} \right)\\ 
         &=    \frac{n+ 2 }{m+ n+ 2 }\cdot \frac{C_{a}^{2} }{C_{a+ b}^{2} }+ \frac{n+ 1 }{m+ n+ 2 }\cdot \frac{C_{a}^{1}C_{b}^{1} }{C_{a+ b}^{2} }+ \frac{n }{m+ n+ 2 }\frac{C_{b}^{2} }{C_{a+ b}^{2} }\\ 
          &= \frac{\left( n+ 2 \right) a\left( a-1 \right)+ 2\left( n+ 1 \right)ab  + nb\left( b-1 \right)  }{\left( m+ n+ 2 \right)\left( a+ b \right)\left( a+ b-1 \right)    }            \\ 
           &= \frac{n }{m+ n+ 2 }+ \frac{2a }{\left( m+ n+ 2 \right)\left( a+ b \right)   }  
        \end{aligned}
        \]
        \begin{remark}
            另一种方法是考虑从乙口袋中取出的白球的来源. 记 \(  A_1  \)为从乙口袋中取出的球最开始就在乙口袋中的事件, 记 \(  A_2  \)为球是从  甲口袋中放到乙口袋中的事件. 则通过对转移的球的颜色情况进行分类讨论, 可以得到 \[
           \begin{aligned}
            P\left( A_{W}|A_2 \right)&=  \frac{1 }{2 }P\left( B_{WB} \right)+ 1\cdot P\left( B_{WW} \right)   + 0\cdot P\left( B_{BB} \right) \\ 
             &= \frac{a }{a+ b } 
           \end{aligned}
            \]此时 \[
            \begin{aligned}
            P\left( A_{W} \right)&= P\left( A_{W}|A_1 \right)P\left( A_1 \right)+ P\left( A_{W}|A_2 \right)P\left( A_2 \right)      \\ 
             &= \frac{n }{m+ n }\cdot \frac{m+ n }{m+ n+ 2 }+ \frac{ a}{a+ b }\cdot \frac{2 }{m+ n+ 2 }    \\ 
              &= \frac{n }{m+ n+ 2 }+ \frac{2a }{\left( m+ n+ 2 \right)\left( a+ b \right)   }  
            \end{aligned}
            \]结果是相同的.
        \end{remark}
    \end{enumerate}
    
\end{proof}
\begin{exercise}{}{}
有 $n$ 个口袋,每个口袋中均有 $a$ 个白球、$b$ 个黑球.从第一个口袋中任取一球放入第二个口袋,再从第二个口袋中任取一球放入第三个口袋,如此下去,从第 $n-1$ 个口袋中任取一球放入第 $n$ 个口袋,最后从第 $n$ 个口袋中任取一球,求此时取到的是白球的概率.
\end{exercise}

\begin{proof}
    设 \(  A_{k,W} , A_{k,B}  \)分别为从第 \(  k  \)个口袋取到  白球和黑球的事件, \(  k= 1,2,\cdots ,n  \) . 则对于 \(  k\ge 2  \),   \[
    \begin{aligned}
    P\left( A_{k,W} \right)&= P\left( A_{k,W}|A_{k-1,B} \right)P\left( A_{k-1,B} \right)+ P\left( A_{k,W}|A_{k-1,W} \right)     P\left( A_{k-1,W} \right)\\ 
     &= \frac{a }{a+ b+ 1 }  \left( 1-P\left( A_{k-1,W} \right)+ \frac{a+ 1 }{a+ b+ 1 }   \right)P\left( A_{k-1,W} \right)\\ 
      &= \frac{P\left( A_{k-1,W} \right) + a }{a+ b+ 1 }   
    \end{aligned}
    \] 从而 \[
    P\left( A_{k,W} \right)-\frac{a }{a+ b }  = \frac{P\left( A_{k-1,W} \right)-\frac{a }{a+ b }   }{ a+ b+ 1} 
    \]我们有 \[
    P\left( A_{n,W} \right)-\frac{a }{a+ b }= \frac{P\left( A_{1},W \right)-\frac{a }{a+ b }   }{ \left( a+ b+ 1 \right)^{n-1} }= 0   
    \]因此 \[
    P\left( A_{n,W} \right)= \frac{a }{a+ b }  
    \]
\end{proof}
\begin{exercise}{}{}
钥匙掉了,掉在宿舍里、教室里、路上的概率分别是 50\%、30\%和 20\%,而掉在上述三处地方被找到的概率分别是 0.8、0.3 和 0.1.试求找到钥匙的概率.
\end{exercise}

\begin{proof}
    设掉在宿舍里, 教室里, 路上的事件分别为 \(  B_1,B_2,B_3  \), 设找到钥匙的事件为 \(  A  \). 则找到钥匙的概率为 \[
   \begin{aligned}
    P\left( A \right)&= P\left( A|B_1 \right)P\left( B_1 \right)+ P\left( A|B_2 \right)P\left( B_2 \right)+ P\left( A|B_3 \right)P\left( B_3 \right) \\ 
     &= 0.8\times 0.5+ 0.3\times 0.3+ 0.1\times 0.2\\ 
      &= 0.51
   \end{aligned}       
    \] 
\end{proof}
\begin{exercise}{}{}
两台车床加工同样的零件,第一台出现不合格品的概率是 0.03,第二台出现不合格品的概率是 0.06,加工出来的零件放在一起,并且已知第一台加工的零件比第二台加工的零件多一倍.
\begin{enumerate}
    \item 求任取一个零件是合格品的概率;
    \item 如果取出的零件是不合格品,求它是由第二台车床加工的概率.
\end{enumerate}
\end{exercise}

\begin{proof}
    记零件由第一台车床加工的事件为 \(  B_1  \), 由第二台加工的事件为 \(  B_2  \), 则 \[
    P\left( B_1 \right)= 2P\left( B_2 \right)\implies P\left( B_1 \right)= \frac{2 }{3 },\quad P\left( B_2 \right)= \frac{1 }{3 }      
    \]  \begin{enumerate}
        \item 计合格的事件为 \(  A_1 \), 则 \[
        \begin{aligned}
        P\left( A_1 \right)&= P\left( A_1|B_1 \right)P\left( B_1 \right)+ P\left( A_1|B_2 \right)P\left( B_2 \right)      \\ 
         &= \left( 1-0.03 \right)\times \frac{2 }{3 }  + \left( 1-0.06 \right)\times \frac{1 }{3 }  \\ 
          &= 0.96
        \end{aligned}
        \] 
        \item 记取出不合格品的事件为 \(  A_2  \) , 则由Bayes公式, 它是第二台车床加工的概率为 \[
        P\left( B_2|A_2 \right)= \frac{P\left( A_2|B_2\right)P\left( B_2 \right)   }{ P\left( A_2 \right) }  = \frac{0.06\times \frac{1 }{3 }  }{1-0.96 }= 0.5 
        \]
    \end{enumerate}
    
\end{proof}
\begin{exercise}{}{}
有两箱零件,第一箱装 50 件,其中 20 件是一等品.第二箱装 30 件,其中 18 件是一等品.现从两箱中随意挑出一箱,然后从该箱中先后任取两个零件,试求:
\begin{enumerate}
    \item 第一次取出的零件是一等品的概率;
    \item 在第一次取出的是一等品的条件下,第二次取出的零件仍然是一等品的概率.
\end{enumerate}
\end{exercise}

\begin{proof}
    设零件来源于第一箱的事件为 \(  B_1  \), 来源于第二箱的为 \(  B_2  \).  \begin{enumerate}
        \item 设第一次取出的零件是一等品的事件为 \(  A_1  \) , 则 \[
        \begin{aligned}
        P\left( A_1 \right)&= P\left( A_1|B_1 \right)P\left( B_1 \right)+ P\left( A_1|B_2 \right)P\left( B_2 \right)\\ 
         &= \frac{20 }{50 }\cdot \frac{1 }{2 }+ \frac{18 }{30 }\cdot \frac{1 }{2 }\\ 
          &= 0.5          
        \end{aligned}
        \]
        \item 设两次均取出一等品的事件为 \(  A  \). 第二次取出一等品的事件为 \(  A_2  \). 则题述事件为 \(  A|A_1  \), 概率为 \[
    \begin{aligned}
      P\left( A|A_1 \right)&= \frac{P\left( A\cap A_1 \right)  }{P\left( A_1 \right)  }= \frac{P\left( A_1\cap A_2 \right)  }{P\left( A_1 \right)  }            
    \end{aligned} 
        \]  其中 \[
       \begin{aligned}
        P\left( A_1\cap A_2 \right) &=P\left( A_1\cap A_2\cap B_1 \right)+ P\left( A_1\cap A_2\cap B_2 \right)\\ 
         &= P\left( A_2|A_1\cap B_1 \right)P\left(A_1| B_1 \right)P\left( B_1 \right) + P\left( A_2|A_1\cap B_2 \right)      P\left( A_1|B_2 \right) P\left( B_2 \right)\\ 
          &= \frac{19 }{49 }\cdot \frac{20 }{50 } \cdot \frac{1 }{2 }+ \frac{17 }{29 }\cdot \frac{18 }{30 } \cdot \frac{1 }{2 }     \\ 
           &= \frac{3601 }{14210 } 
       \end{aligned}
        \] 因此 \[
        \begin{aligned}
       P\left( A|A_1 \right)= \frac{P\left( A_1\cap A_2 \right)  }{P\left( A_1 \right)  }= \frac{3601 }{7105 }   
        \end{aligned}   
        \]
    \end{enumerate}
    
\end{proof}
\begin{exercise}{}{}
学生在做一道有 4 个选项的单项选择题时,如果他不知道问题的正确答案,就作随机猜测.现从卷面上看题是答对了,试在以下情况下求学生确实知道正确答案的概率:
\begin{enumerate}
    \item 学生知道正确答案和胡乱猜测的概率都是 1/2;
    \item 学生知道正确答案的概率是 0.2.
\end{enumerate}
\end{exercise}

\begin{proof}
   设学生知道正确答案的事件为 \(  A_1  \), 不知道正确答案的 事件为 \(  A_2  \), 设学生卷面上答对的事件为 \(  B_1  \), 答错的事件为 \(  B_2  \).题述事件为 \(  A_1|B_1  \)  
   \begin{enumerate}
    \item 当 \(  P\left( A_1 \right)= P\left( A_2 \right)= \frac{1 }{2 }     \)时, 题述事件概率为 \[
    \begin{aligned}
    P\left( A_1|B_1 \right)&= \frac{P\left( B_1|A_1 \right)P\left( A_1 \right)   }{P\left( B_1 \right)  }   \\ 
     &= \frac{P\left( B_1|A_1 \right)P\left( A_1 \right)   }{P\left( B_1|A_1 \right)P\left( A_1 \right)+ P\left( B_1|A_2 \right)P\left( A_2 \right)     } \\ 
      &= \frac{1\cdot \frac{1 }{2 }  }{1\cdot \frac{1 }{2 }+ \frac{1 }{4 }\cdot \frac{1 }{2 }    }\\ 
       &= 0.8 
    \end{aligned}
    \]
    \item 当 \(  P\left( A_1 \right)= 0.2   \)时 , 题述事件概率为 \[
  \begin{aligned}
    P\left( A_1|B_1 \right)&= \frac{P\left( B_1|A_1 \right)P\left( A_1 \right)   }{P\left( B_1|A_1 \right)P\left( A_1 \right)+ P\left( B_1|A_2 \right)P\left( A_2 \right)     } \\ 
      &= \frac{1\cdot 0.2 }{1\cdot 0.2+ \frac{1 }{4 }\cdot 0.8  }\\ 
       &= 0.5  
  \end{aligned}
    \] 
   \end{enumerate}
      
\end{proof}


\begin{exercise}{}{}
口袋中有一个球,不知它的颜色是黑的还是白的.现再往口袋中放入一个白球,然后从口袋中任意取出一个,发现取出的是白球,试问口袋中原来那个球是白球的可能性为多少?
\end{exercise}
\begin{proof}
    设若口袋中随机出现一个球, 为白球的事件为 \(  A_1  \), 为黑球的事件为 \(  A_2  \) . 设后从口袋中取出的是白球的事件为 \(  B_1  \), 取出黑球的事件为 \(  B_2  \) . 则题述事件为 \(  A_1|B_1  \), 概率为     \[
   \begin{aligned}
    P\left( A_1|B_1 \right) &= \frac{P\left( B_1|A_1 \right)P\left( A_1 \right)   }{P\left( B_1|A_1 \right)P\left( A_1 \right)+ P\left( B_1|A_2 \right)P\left( A_2 \right)     }  \\ 
     &= \frac{1\cdot \frac{1 }{2 }  }{1\cdot \frac{1 }{2 }+ \frac{1 }{2 }\cdot \frac{1 }{2 }    }\\ 
      &= \frac{2 }{3 }  
   \end{aligned}
    \]
\end{proof}
\begin{exercise}{}{}
$m$ 个人相互传球, 球从甲手中开始传出, 每次传球时, 传球者等可能地把球传给其余 $m-1$ 个人中的任何一个. 求第 $n$ 次传球时仍由甲传出的概率.
\end{exercise}
\begin{proof}
    设第 \(  k  \)次传球时由甲传出的概率为 \(  A_{k}  \). 则 对于所有的 \(  k\ge 2  \), \[
  \begin{aligned}
    P\left( A_{k} \right)&= P\left( A_{k}|A_{k-1} \right)P\left( A_{k-1} \right)+ P\left( A_{k} |A_{k-1}^{c}\right)P\left( A_{k-1}^{c} \right) \\
     &=0+ \frac{1 }{m-1 }P\left( A_{k-1}^{c} \right)\\ 
      &= \frac{1-P\left( A_{k-1} \right)  }{m-1 }   
  \end{aligned}     
    \]从而 \[
    P\left( A_{k} \right)-\frac{1 }{m }= \frac{-\left( P\left( A_{k-1} \right)-\frac{1 }{m }   \right)  }{m-1 }   
    \]得到 \[
    P\left( A_{n} \right)-\frac{1 }{m }= \left( -\frac{1 }{m-1 }  \right)^{n-1}\left( 1-\frac{1 }{m }   \right)= \left( -\frac{1 }{m-1 }  \right)^{n-1}\left( \frac{m-1 }{m }  \right)      
    \]因此 \[
    P\left( A_{n} \right) = \frac{1 }{m }\left( 1+  \frac{\left( -1 \right)^{n-1}  }{\left( m-1 \right)^{n-2}  }  \right) (n\ge 2),\quad P\left( A_1 \right)= 1   
    \]
\end{proof}
\begin{exercise}{}{}
甲口袋有 1 个黑球、2 个白球, 乙口袋有 3 个白球. 每次从两口袋中各任取一球, 交换后放入另一口袋. 求交换 $n$ 次后, 黑球仍在甲口袋中的概率.
\end{exercise}

\begin{proof}
    设 \(  A_{k}  \)表示交换 \(  k  \)次后黑球仍在甲口袋中的事件. 则 \[
   \begin{aligned}
    P\left( A_{k} \right)&= P\left( A_{k}|A_{k-1} \right)P\left( A_{k-1} \right)+ P\left( A_{k}|A_{k-1}^{c} \right)P\left( A_{k-1}^{c} \right)\\ 
     &=  \frac{2 }{3 }\cdot P\left( A_{k-1} \right)  + \frac{1 }{3 }\cdot P\left( A_{k-1}^{c} \right)  \\ 
      &= \frac{2 }{3 }P\left( A_{k-1} \right)+ \frac{1 }{3 }\left( 1-P\left( A_{k-1} \right)  \right)    \\ 
       &= \frac{1 }{3 }P\left( A_{k-1} \right)+ \frac{1 }{3 }   
   \end{aligned}     
    \]  注意到 \[
    P\left( A_{k} \right)-\frac{1 }{2 } =  \frac{1 }{3 }\left( P\left( A_{k-1} \right)-\frac{1 }{2}   \right)    
    \]我们有 \[
    P\left( A_{n} \right)-\frac{1 }{2 }= \frac{1 }{3^{n-1} }\left( P\left( A_{1} \right)-\frac{1 }{2 }   \right)    
    \]其中 \(  P\left( A_1 \right)= \frac{2 }{3 }    \) , 带入得到 \[
    P\left( A_{n} \right)=\frac{1 }{2 }\left( \frac{1 }{3^{n} }+ 1  \right)  
    \]
\end{proof}
\begin{exercise}{}{}
口袋中有 $a$ 个白球, $b$ 个黑球和 $n$ 个红球, 现从中一个一个不返回地取球. 试证白球比黑球出现得早的概率为 $a/(a+b)$, 与 $n$ 无关.
\end{exercise}

\begin{proof}
    设 \(  B_{1}, W_{1},R_{1}  \)为第 \(  1  \)次取到黑, 白, 红球的事件. 记 \(  A_{k}  \)为有 \(  k  \)个红球时, 白球比黑球出现得早的概率. 则对于所有的 \(  k\ge 2  \)  \[
   \begin{aligned}
    P\left( A_{k} \right)&= P\left( A_{k}|R_{1} \right)P\left( R_1 \right)+ P\left( A_{k}|W_{1} \right)P\left( W_{1} \right)+ P\left( A_{k}|B_1 \right)P\left( B_1 \right)\\ 
     &=        P\left( A_{k-1} \right)\cdot \frac{k }{a+ b+ k }+ 1\cdot \frac{a }{a+ b + k}+ 0\\ 
      &=     \frac{P\left( A_{k-1} \right)k+ a  }{a+ b+ k } 
   \end{aligned}
    \]   \[
  \begin{aligned}
    P\left( A_{k} \right)-\frac{a }{a+ b }&= \frac{P\left( A_{k-1} \right)k+ a  }{a+ b+ k}-\frac{a }{a+ b }   \\ 
     &= \frac{P\left( A_{k-1} \right)k\left( a+ b \right)+ a\left( a+ b \right)-a\left( a+ b+ k \right)     }{ \left( a+ b+ k \right)\left( a+ b \right)  }   \\ 
      &= \frac{k\left( P\left( A_{k-1} \right)\left( a+ b \right)-a   \right)  }{\left( a+ b+ k \right)\left( a+ b \right)   }\\ 
       &= \frac{k\left( P\left( A_{k-1} \right)-\frac{a }{a+ b }   \right)  }{\left( a+ b+ k \right)  }  
  \end{aligned}
    \]易见 \(  P\left( A_1 \right)= \frac{a }{a+ b }    \), 且当 \(  P\left( A_{k-1} \right)= \frac{a }{a+ b }    \)  时,  \(  P\left( A_{k} \right)  = \frac{a }{a+ b }  \), 归纳可知 \[
    P\left( A_{n} \right)= \frac{a }{a+ b }  
    \] 
\end{proof}
\end{document}